% Sección 6
\section{Qué mide hoy el prototipo}\label{sec:prototipo}

\subsection{Los indicadores del reto}

La ficha de la vertical 04 mide el resultado con dos indicadores. El primero es el porcentaje de entregas a tiempo en puntos críticos, como un albergue o un comedor comunitario. El segundo es el desperdicio de alimentos perecederos, que debe quedar por debajo del 5~\%. Hoy los medimos solo en simulación, porque el prototipo no ha operado con datos reales de Manizales.

\subsection{Cómo se simuló}

La simulación compara cuatro formas de despachar la misma jornada, con los mismos pedidos, la misma flota y los mismos cierres. Usa el grafo de conducción de Manizales y Villamaría de OpenStreetMap. Cada escenario tiene entre 30 y 80 puntos de entrega, de los cuales el 20~\% son críticos. La flota mezcla camiones, camperos y motos, y en cada jornada se cierran entre 3 y 8 tramos de vía. Nadie conoce un cierre hasta que un vehículo intenta pasar por él. La jornada dura 14 horas, y lo que no se entregó para entonces cuenta como no entregado.

Las cuatro políticas se distinguen por cómo deciden. La línea base, llamada ingenua, atiende por orden de llegada y saca los lotes de bodega por orden de ingreso. La política experta imita a un despachador con oficio, que atiende primero lo urgente y saca primero lo que vence. Las otras dos usan el optimizador de DespachoX. Una lo llama una sola vez al inicio, y la otra, que es DespachoX completo, recalcula con cada cierre, cada tanda de pedidos y cada regreso de un vehículo al acopio.

Cada nivel de carga se corrió con 30 semillas, es decir, 30 jornadas sintéticas, las mismas para todas las políticas. Los resultados se comparan por pares, semilla a semilla, con un intervalo de confianza del 95~\%. Todos los datos de demanda, flota, inventario y cierres son sintéticos. Las velocidades, los tiempos y los plazos son supuestos del equipo, y ninguno viene de una operación real del Banco de Alimentos ni de los organismos de socorro. Nadie del equipo ha revisado todavía estos resultados.

\subsection{Tres cifras con su línea base}

La tabla~\ref{tab:simulacion} trae las tres cifras que el agente de simulación propuso para el pitch, con su línea base y su número de escenarios. Todas son simulación con supuestos nuestros, y su nivel de evidencia es \textbf{Inferencia}. Las fuentes son las entradas de las 4:57~p.\,m. del canal del equipo (\ruta{torre/CANAL.md}) y la sección 8 de \ruta{docs/SIMULACION.md}, en la rama principal del repositorio. La corrida con flota corta todavía no está descrita en \ruta{docs/SIMULACION.md}, así que su número de semillas sale de su archivo de resultados.

{\footnotesize
\setlength{\tabcolsep}{4pt}
\renewcommand{\arraystretch}{1.3}
\begin{longtable}{L{0.20\linewidth} L{0.15\linewidth} L{0.24\linewidth} L{0.32\linewidth}}
\caption{Cifras de simulación con su línea base. Ninguna es una medición en la operación real.}\label{tab:simulacion}\\
\toprule
\textbf{Indicador} & \textbf{DespachoX} & \textbf{Línea base} & \textbf{Escenarios y condiciones}\\
\midrule
\endfirsthead
\multicolumn{4}{l}{\small\itshape Tabla~\ref{tab:simulacion}, continuación}\\
\toprule
\textbf{Indicador} & \textbf{DespachoX} & \textbf{Línea base} & \textbf{Escenarios y condiciones}\\
\midrule
\endhead
\midrule
\multicolumn{4}{r}{\small\itshape Sigue en la página siguiente}\\
\endfoot
\bottomrule
\endlastfoot
(1) Entregas a tiempo en puntos críticos, carga media & 99,5~\% & 70,9~\% con despacho por orden de llegada. Diferencia de +28,6 puntos, con intervalo [22,4; 34,7] & Simulación con 30 semillas nuevas que no intervinieron en ninguna decisión. Usa un peso de 3 para los puntos críticos (\texttt{FACTOR\_CRITICO}), propuesto a las 4:57~p.\,m. y que a esa hora no era todavía el valor por defecto del código.\\
(2) Perecederos vencidos, carga media & 0,0~\% & 1,8~\% al sacar de bodega por orden de ingreso & Simulación con 30 semillas. Con dos camperos menos, también en 30 semillas, el orden de ingreso llega al 5,7~\% y pasa la meta del 5~\%, mientras DespachoX sigue en 0,0~\%.\\
(3) Entregas a tiempo en todos los puntos, con un 10~\% de pedidos mal ubicados & +6,0 puntos, con intervalo [4,5; 7,5] & El mismo sistema sin el paso de confirmar la ubicación con quien pidió & Simulación con 30 semillas, carga media. El 10~\% es un supuesto sin fuente, porque no sabemos cuántos pedidos llegan mal ubicados.\\
\end{longtable}
}

La primera cifra cambia poco con ese peso. Con el peso de 1 que usa hoy el código, la primera ronda dio una diferencia parecida en carga media, de 28,9 puntos con intervalo [23,2; 34,7]. La ventaja frente al despacho por orden de llegada crece con la carga. En la primera ronda fue de 12, 29 y 46 puntos en carga holgada, media y alta, y no sabemos cuál de esas cargas se parece a la real.

\subsection{Lo que hay que saber decir frente a un despachador experto}

Un despachador con experiencia consigue casi lo mismo que DespachoX en los dos indicadores del reto, y conviene decirlo antes de que lo diga el jurado. En entregas a tiempo en puntos críticos, DespachoX empata con la política experta. En las 30 semillas nuevas, con el peso de 3, la diferencia fue de 0,3, 1,7 y 0,8 puntos según la carga, y todos los intervalos incluyen el cero. En entregas a tiempo de todos los puntos, DespachoX queda unos 2 puntos por debajo en carga media y alta. En perecederos, los dos quedan en 0,0~\%, porque el desperdicio lo resuelve una regla de bodega y no el ruteo.

Hay dos situaciones simuladas en que DespachoX supera a la experta. Con la flota corta, de dos camperos menos, en carga media y con 30 semillas, DespachoX entrega a tiempo 3,4 puntos más en todos los puntos, con intervalo [1,3; 5,5]. Con un 10~\% de pedidos mal ubicados, también en carga media y con 30 semillas, confirmar la ubicación le da 3,0 puntos más en puntos críticos [0,9; 5,2] y 3,4 más en todos los puntos [2,1; 4,8]. Esta segunda ventaja depende de un supuesto, que el despachador experto no llama a confirmar cada pedido. Si lo hiciera, tendría el mismo beneficio a costa de su tiempo, que la simulación no modela.

DespachoX también recorre menos kilómetros. En las semillas nuevas, su flota hizo entre 47 y 60 km menos por jornada que la política experta. En la primera ronda la diferencia equivalía a entre el 5~\% y el 11~\% del recorrido.

La prioridad tiene un costo en equidad. Atender primero lo crítico deja para el final a los puntos de baja urgencia. En la primera ronda, con DespachoX el punto peor atendido esperó 2,3 horas más en carga media y 4,3 horas más en carga alta que con el despacho por orden de llegada.

Los cierres viales pesan poco en estos resultados. Con 0, 5 o 10 cierres por jornada, la ventaja sobre la línea base casi no cambia, porque viene de priorizar los pedidos críticos. La red urbana simulada tiene caminos alternos, y el aviso por radio desvía a todos los conductores por igual. Esto puede cambiar en la red de la región, donde hay menos alternativas. Esa última idea es una hipótesis que la simulación todavía no prueba.

Frente a un experto, entonces, el argumento de DespachoX es la disponibilidad de su criterio. El sistema pone las reglas de un despachador con oficio al alcance de cualquier voluntario, a cualquier hora. Además las combina con información que hoy llega tarde o no llega, como los pedidos del chat, los cierres que reportan los conductores y la confirmación de cada entrega. La simulación todavía no mide el valor de esa información, salvo en el experimento de la ubicación.

\subsection{Lo que la simulación no prueba}

\begin{itemize}
\item No prueba que DespachoX supere a un buen despachador humano en los indicadores del reto.
\item No mide el valor de recibir los pedidos por chat ni el de la trazabilidad, salvo el paso de confirmar la ubicación.
\item No valida sus supuestos. Las velocidades, los tiempos de carga y de entrega, los plazos, la flota, la demanda y los cierres son nuestros.
\item No representa bien los perecederos. Hay un solo acopio, tanto inventario perecedero como demanda, y lotes que vencen entre 6 y 48 horas.
\item No incluye errores humanos ni del modelo de lenguaje, y cubre una sola jornada de 14 horas.
\end{itemize}

Según el canal, a las 4:57~p.\,m. seguían en curso dos corridas. Una usa 3 segundos de búsqueda por llamada al optimizador, y la otra corre la carga media sobre la red ampliada de la región. El agente de simulación esperaba tenerlas hacia las 5:45~p.\,m.
